\section{The higher-rung \texorpdfstring{\(\mathrm{SU}(3)\)}{SU(3)} Gram cone}
\label{sec:higher-rung-su3-gram-cone}

This section isolates the part of the proposed higher-rung construction that
is independent of the candidate complex threefold and proves it from the
underlying linear action.  The compact object used here is the representation
sphere
\[
 \operatorname{Mat}_{3\times2}(\mathbb C)^+,
\]
the one-point compactification of a twelve-real-dimensional vector space.
It must not be identified with a cusp-induced \(S^{12}\) construction unless
an additional based identification has first been chosen.  The proofs below
also distinguish three structures that otherwise look deceptively similar:
the Lorentzian determinant on \(\operatorname{Herm}_2(\mathbb C)\), the
positive-semidefinite Gram cone inside that vector space, and a separately
chosen Euclidean metric on the positive-definite cone.

\subsection{Conventions and the Gram morphism}

Put
\begin{equation}
 \mathscr V=\operatorname{Mat}_{3\times2}(\mathbb C),
 \qquad
 \Gamma:\mathscr V\longrightarrow\operatorname{Herm}_2(\mathbb C),
 \qquad
 \Gamma(A)=A^\dagger A.
 \label{eq:hrgc-gram-map}
\end{equation}
Write \(A=(a\ b)\), with \(a,b\in\mathbb C^3\), and let
\(\mathrm{SU}(3)\) act on the left:
\[
 U\cdot A=UA.
\]
For \(r=0,1,2\), let
\[
 \mathscr V_{(r)}=\{A\in\mathscr V:\operatorname{rank}_{\mathbb C}A=r\}.
\]
The Hermitian products are linear in the second argument.  We orient
\(\mathscr V\) by the complex orientation, in the real coordinate order
\begin{equation}
 \begin{split}
 (&\operatorname{Re}A_{11},\operatorname{Im}A_{11},
   \operatorname{Re}A_{12},\operatorname{Im}A_{12},
   \operatorname{Re}A_{21},\operatorname{Im}A_{21},\\
  &\operatorname{Re}A_{22},\operatorname{Im}A_{22},
   \operatorname{Re}A_{31},\operatorname{Im}A_{31},
   \operatorname{Re}A_{32},\operatorname{Im}A_{32}).
 \end{split}
 \label{eq:hrgc-v-orientation}
\end{equation}
On \(\operatorname{Herm}_2(\mathbb C)\), use the oriented real-linear
coordinates
\begin{equation}
 \Phi(t,x,y,z)=
 \begin{pmatrix}
  t+z&x-iy\\
  x+iy&t-z
 \end{pmatrix},
 \qquad
 \operatorname{or}_{\operatorname{Herm}}=dt\wedge dx\wedge dy\wedge dz.
 \label{eq:hrgc-pauli-coordinates}
\end{equation}
Thus, if \(\Gamma(A)_{12}=\langle a,b\rangle\), then
\begin{equation}
 \begin{aligned}
 t&=\frac12(\lVert a\rVert^2+\lVert b\rVert^2),&
 z&=\frac12(\lVert a\rVert^2-\lVert b\rVert^2),\\
 x&=\operatorname{Re}\langle a,b\rangle,&
 y&=-\operatorname{Im}\langle a,b\rangle.
 \end{aligned}
 \label{eq:hrgc-column-pauli-coordinates}
\end{equation}

\subsection{Exact orbit classification}

\begin{theorem}[The Gram quotient]
\label{thm:hrgc-gram-quotient}
The morphism \(\Gamma\) is \(\mathrm{SU}(3)\)-invariant, has image exactly
\(\operatorname{Herm}_2(\mathbb C)_{\geq0}\), and separates
\(\mathrm{SU}(3)\)-orbits.  It induces a proper homeomorphism
\begin{equation}
 \overline\Gamma:
 \mathscr V/\mathrm{SU}(3)
 \xrightarrow{\ \cong\ }
 \operatorname{Herm}_2(\mathbb C)_{\geq0}.
 \label{eq:hrgc-orbit-homeomorphism}
\end{equation}
\end{theorem}

\begin{proof}
For \(U\in\mathrm{SU}(3)\),
\[
 \Gamma(UA)=A^\dagger U^\dagger UA=A^\dagger A,
\]
so \(\Gamma\) is invariant.  Every matrix \(A^\dagger A\) is positive
semidefinite.  Conversely, for
\(H\in\operatorname{Herm}_2(\mathbb C)_{\geq0}\), let \(H^{1/2}\) be its
unique positive-semidefinite square root and define
\begin{equation}
 s(H)=
 \begin{pmatrix}
  H^{1/2}\\
  0\quad0
 \end{pmatrix}\in\operatorname{Mat}_{3\times2}(\mathbb C).
 \label{eq:hrgc-square-root-section}
\end{equation}
Then \(s(H)^\dagger s(H)=H\), proving surjectivity and giving a continuous
section.

It remains to prove exact orbit separation.  Suppose
\(A^\dagger A=B^\dagger B\).  Define
\[
 T:\operatorname{im}A\longrightarrow\operatorname{im}B,
 \qquad T(Au)=Bu.
\]
If \(Au=0\), then
\[
 0=\lVert Au\rVert^2
  =u^\dagger A^\dagger Au
  =u^\dagger B^\dagger Bu
  =\lVert Bu\rVert^2,
\]
so \(Bu=0\), and \(T\) is well defined.  For \(u,v\in\mathbb C^2\),
\[
 \langle T(Au),T(Av)\rangle
 =u^\dagger B^\dagger Bv
 =u^\dagger A^\dagger Av
 =\langle Au,Av\rangle.
\]
Thus \(T\) is a unitary isomorphism of the image subspaces.  Extend it to
some \(U_0\in\mathrm U(3)\) satisfying \(U_0A=B\).  Since
\(\operatorname{rank}B\leq2\), the orthogonal complement of
\(\operatorname{im}B\) has positive complex dimension.  Choose a unitary
map \(V\) which fixes \(\operatorname{im}B\) pointwise and whose determinant
is \((\det U_0)^{-1}\).  Then
\[
 U=VU_0\in\mathrm{SU}(3),
 \qquad UA=VB=B.
\]
Hence equal Gram matrices are exactly equal \(\mathrm{SU}(3)\)-orbits.

The action is proper because \(\mathrm{SU}(3)\) is compact; hence the orbit
space is Hausdorff.  The continuous map
\(H\mapsto[s(H)]\) is inverse to the continuous induced map
\(\overline\Gamma\), so \eqref{eq:hrgc-orbit-homeomorphism} is a
homeomorphism.  Finally,
\begin{equation}
 \lVert A\rVert_F^2=\operatorname{tr}(A^\dagger A)
 =\operatorname{tr}\Gamma(A).
 \label{eq:hrgc-properness-identity}
\end{equation}
If \(K\) is compact in the positive-semidefinite cone, then
\(\Gamma^{-1}(K)\) is closed and bounded by
\eqref{eq:hrgc-properness-identity}, hence compact.  Thus \(\Gamma\) is
proper.  For compact \(K\),
\(\overline\Gamma^{-1}(K)\) is the continuous image in the orbit space of
the compact set \(\Gamma^{-1}(K)\), so \(\overline\Gamma\) is proper as
well.
\end{proof}

\begin{proposition}[Stabilizers and orbit-type strata]
\label{prop:hrgc-stabilizers}
For \(A\in\mathscr V_{(r)}\),
\begin{equation}
 \operatorname{Stab}_{\mathrm{SU}(3)}(A)\cong\mathrm{SU}(3-r).
 \label{eq:hrgc-stabilizer-formula}
\end{equation}
The exact dimensions are
\begin{equation}
 \begin{array}{c|c|c|c|c}
 r&\dim_{\mathbb R}\mathscr V_{(r)}&\operatorname{Stab}(A)&
 \dim_{\mathbb R}(\mathrm{SU}(3)\cdot A)&
 \dim_{\mathbb R}(\mathscr V_{(r)}/\mathrm{SU}(3))\\ \hline
 2&12&\{1\}&8&4\\
 1&8&\mathrm{SU}(2)&5&3\\
 0&0&\mathrm{SU}(3)&0&0.
 \end{array}
 \label{eq:hrgc-stratum-table}
\end{equation}
Consequently the quotient has topological dimension four, but the action is
free only on the rank-two stratum.
\end{proposition}

\begin{proof}
The equality \(UA=A\) is equivalent to the assertion that \(U\) fixes
\(\operatorname{im}A\) pointwise.  Restriction to its orthogonal complement
identifies the stabilizer with \(\mathrm{SU}(3-r)\), proving
\eqref{eq:hrgc-stabilizer-formula}.  The rank-two locus is open in the
twelve-real-dimensional vector space.  The complex rank-one locus in the
space of \(3\times2\) matrices has complex dimension
\[
 1(3+2-1)=4,
\]
hence real dimension eight away from the origin.  Subtracting
\(8-\dim_{\mathbb R}\operatorname{Stab}(A)\) gives the displayed orbit and
quotient dimensions.  They agree with the target strata: positive-definite
Hermitian matrices form a four-dimensional open cone, while every nonzero
rank-one positive-semidefinite matrix has a unique expression
\(\lambda vv^\dagger\) with \(\lambda>0\) and \(v\in S^3\), modulo
\(v\sim e^{i\theta}v\).  Hence that target stratum is
\[
 (0,\infty)\times\mathbb{CP}^1
 \cong(0,\infty)\times S^2,
\]
of real dimension three.  The rank-zero stratum is the origin.
\end{proof}

\subsection{The free locus, its explicit trivialization, and orientation}

\begin{theorem}[Determinant-one completion]
\label{thm:hrgc-free-trivialization}
There is an \(\mathrm{SU}(3)\)-equivariant diffeomorphism
\begin{equation}
 \mathscr V_{(2)}
 \xrightarrow{\ \cong\ }
 \mathrm{SU}(3)\times\operatorname{Herm}_2(\mathbb C)_{>0}
 \label{eq:hrgc-free-product}
\end{equation}
under which \(\Gamma\) is projection to the second factor and the group acts
by left multiplication on the first.  Thus the principal bundle over the
free quotient has a global section and is trivial.
\end{theorem}

\begin{proof}
For \(A\in\mathscr V_{(2)}\), set
\[
 H=A^\dagger A>0,
 \qquad Q=AH^{-1/2}=(q_1\ q_2).
\]
Then \(Q^\dagger Q=I_2\), so \((q_1,q_2)\) is an ordered orthonormal
two-frame in \(\mathbb C^3\).  With \(\varepsilon_{123}=1\), define
\begin{equation}
 (q_3)_i=\sum_{j,k=1}^3\varepsilon_{ijk}
 \overline{(q_1)_j}\,\overline{(q_2)_k}.
 \label{eq:hrgc-cross-product}
\end{equation}
Antisymmetry gives
\(q_1^\dagger q_3=q_2^\dagger q_3=0\).  The Gram--Pl\"ucker identity, proved
independently in \cref{prop:hrgc-pluecker}, gives \(\lVert q_3\rVert=1\).
Finally, contraction of the two alternating tensors gives
\begin{align*}
 \det[q_1\ q_2\ q_3]
 &=\sum_{i,j,k}\varepsilon_{ijk}(q_1)_i(q_2)_j(q_3)_k\\
 &=\lVert q_1\rVert^2\lVert q_2\rVert^2
   -|q_1^\dagger q_2|^2=1.
\end{align*}
Thus \(U_A=[q_1\ q_2\ q_3]\in\mathrm{SU}(3)\), and
\[
 A=U_A s(H),
\]
where \(s\) is \eqref{eq:hrgc-square-root-section}.  The positive square
root, inverse square root, and the polynomial-conjugate formula
\eqref{eq:hrgc-cross-product} are smooth on the positive-definite locus.
Consequently
\[
 A\longmapsto(U_A,A^\dagger A),
 \qquad
 (U,H)\longmapsto Us(H)
\]
are smooth inverse maps.  Their equivariance is immediate.
\end{proof}

The order of factors in \eqref{eq:hrgc-free-product} does not by itself fix
an orientation.  We now fix and calculate it.  Give
\(\mathfrak{su}(3)\) the ordered real basis
\begin{equation}
 \begin{gathered}
 X_{12},Y_{12},X_{13},Y_{13},X_{23},Y_{23},D_1,D_2,\\
 X_{jk}=E_{jk}-E_{kj},\qquad
 Y_{jk}=i(E_{jk}+E_{kj}),\\
 D_1=i\operatorname{diag}(1,-1,0),\qquad
 D_2=i\operatorname{diag}(0,1,-1).
 \end{gathered}
 \label{eq:hrgc-su3-oriented-basis}
\end{equation}
At
\[
 A_0=\begin{pmatrix}1&0\\0&1\\0&0\end{pmatrix},
\]
use the oriented Hermitian basis
\begin{equation}
 K_t=I_2,\quad
 K_x=\begin{pmatrix}0&1\\1&0\end{pmatrix},\quad
 K_y=\begin{pmatrix}0&-i\\i&0\end{pmatrix},\quad
 K_z=\begin{pmatrix}1&0\\0&-1\end{pmatrix}
 \label{eq:hrgc-hermitian-oriented-basis}
\end{equation}
and horizontal lifts \(B_\alpha=\tfrac12A_0K_\alpha\).  Directly,
\(d\Gamma_{A_0}(B_\alpha)=K_\alpha\).  Relative to
\eqref{eq:hrgc-v-orientation}, the twelve columns
\[
 (X_{12}A_0,Y_{12}A_0,X_{13}A_0,Y_{13}A_0,
 X_{23}A_0,Y_{23}A_0,D_1A_0,D_2A_0,B_t,B_x,B_y,B_z)
\]
form the exact real coordinate matrix
\begingroup
\scriptsize
\setlength{\arraycolsep}{2.4pt}
\begin{equation}
 M_{\mathrm{or}}:=\left(
 \begin{array}{rrrrrrrrrrrr}
0&0&0&0&0&0&0&0&\frac12&0&0&\frac12\\
0&0&0&0&0&0&1&0&0&0&0&0\\
1&0&0&0&0&0&0&0&0&\frac12&0&0\\
0&1&0&0&0&0&0&0&0&0&-\frac12&0\\
-1&0&0&0&0&0&0&0&0&\frac12&0&0\\
0&1&0&0&0&0&0&0&0&0&\frac12&0\\
0&0&0&0&0&0&0&0&\frac12&0&0&-\frac12\\
0&0&0&0&0&0&-1&1&0&0&0&0\\
0&0&-1&0&0&0&0&0&0&0&0&0\\
0&0&0&1&0&0&0&0&0&0&0&0\\
0&0&0&0&-1&0&0&0&0&0&0&0\\
0&0&0&0&0&1&0&0&0&0&0&0
 \end{array}
 \right).
 \label{eq:hrgc-orientation-matrix}
\end{equation}
\endgroup
Laplace expansion gives
\begin{equation}
 \det M_{\mathrm{or}}=\frac12>0.
 \label{eq:hrgc-orientation-determinant}
\end{equation}
Since both factors in \eqref{eq:hrgc-free-product} are connected, the sign is
constant on the free locus.  Therefore the quotient orientation determined
by
\[
 \operatorname{or}_{\mathscr V}
 =\operatorname{or}_{\mathrm{SU}(3)}\wedge
  \operatorname{or}_{\mathscr V_{(2)}/\mathrm{SU}(3)}
\]
agrees under \(\Gamma\) with
\(dt\wedge dx\wedge dy\wedge dz\).

\subsection{Lorentz determinant and Pl\"ucker norm}

For \(H=\Phi(t,x,y,z)\), direct expansion gives
\begin{equation}
 \operatorname{tr}H=2t,
 \qquad
 \det H=t^2-x^2-y^2-z^2,
 \label{eq:hrgc-lorentz-determinant}
\end{equation}
and the eigenvalues are
\(t\pm\sqrt{x^2+y^2+z^2}\).  Hence
\begin{equation}
 \begin{aligned}
 H>0&\Longleftrightarrow t>\sqrt{x^2+y^2+z^2},\\
 H\geq0&\Longleftrightarrow t\geq\sqrt{x^2+y^2+z^2},\\
 \operatorname{rank}H=1&\Longleftrightarrow
 t=\sqrt{x^2+y^2+z^2}>0.
 \end{aligned}
 \label{eq:hrgc-positive-cone-inequalities}
\end{equation}
Thus \(\det\) has signature \((1,3)\) when positive directions are listed
first, while \(-\det\) has signature \((3,1)\).  The quotient image is not
all of this Lorentz vector space: it is precisely the closed future causal
cone, and its free locus is the open future timelike cone.

\begin{proposition}[Gram--Pl\"ucker identity]
\label{prop:hrgc-pluecker}
For \(A=(a\ b)\),
\begin{equation}
 \det\Gamma(A)
 =\lVert a\rVert^2\lVert b\rVert^2-|\langle a,b\rangle|^2
 =\sum_{1\leq i<j\leq3}|a_ib_j-a_jb_i|^2
 =\lVert a\wedge b\rVert^2.
 \label{eq:hrgc-pluecker-identity}
\end{equation}
\end{proposition}

\begin{proof}
Expanding gives
\begin{align*}
 \lVert a\rVert^2\lVert b\rVert^2-|\langle a,b\rangle|^2
 &=\sum_{i,j}|a_i|^2|b_j|^2
   -\sum_{i,j}\overline{a_i}b_i a_j\overline{b_j}\\
 &=\sum_{i<j}
 \bigl(|a_i|^2|b_j|^2+|a_j|^2|b_i|^2
 -\overline{a_i}b_i a_j\overline{b_j}
 -\overline{a_j}b_j a_i\overline{b_i}\bigr)\\
 &=\sum_{i<j}|a_ib_j-a_jb_i|^2.
\end{align*}
The last expression is the squared Hermitian norm of \(a\wedge b\) in the
ordered basis
\((e_1\wedge e_2,e_1\wedge e_3,e_2\wedge e_3)\).
\end{proof}

The determinant is only one invariant and does not classify orbits.  For
example,
\[
 A_1=\begin{pmatrix}1&0\\0&1\\0&0\end{pmatrix},
 \qquad
 A_2=\begin{pmatrix}\sqrt2&0\\0&1/\sqrt2\\0&0\end{pmatrix}
\]
satisfy \(\det\Gamma(A_1)=\det\Gamma(A_2)=1\), but their Gram matrices are
\(I_2\) and \(\operatorname{diag}(2,1/2)\), so
\cref{thm:hrgc-gram-quotient} shows that they lie in distinct orbits.

\subsection{The four-coefficient form is additional input}

Let \(r=(r_1,r_2,r_3,r_4)^{\mathsf T}\in\mathbb R^4\), and define the
orthogonal Hadamard coordinates
\begin{equation}
 \begin{pmatrix}h\\x\\y\\z\end{pmatrix}
 =\frac12
 \begin{pmatrix}
 1&1&1&1\\
 1&1&-1&-1\\
 1&-1&1&-1\\
 1&-1&-1&1
 \end{pmatrix}
 \begin{pmatrix}r_1\\r_2\\r_3\\r_4\end{pmatrix}.
 \label{eq:hrgc-hadamard-transform}
\end{equation}
The proposed coefficient-space form
\begin{equation}
 Q_4(r)=\frac12\left(\sum_{i=1}^4r_i\right)^2-
 \sum_{i=1}^4r_i^2
 \label{eq:hrgc-q4-definition}
\end{equation}
satisfies the exact identity
\begin{equation}
 Q_4(r)=h^2-x^2-y^2-z^2.
 \label{eq:hrgc-q4-lorentz}
\end{equation}
Indeed, orthogonality gives
\(\sum r_i^2=h^2+x^2+y^2+z^2\) and
\(\sum r_i=2h\).

The same Hermitian matrix can be written in the alternative convention
\begin{equation}
 \Psi(h,x,y,z)=
 \begin{pmatrix}h+y&x+iz\\x-iz&h-y\end{pmatrix}.
 \label{eq:hrgc-pro-hermitian-convention}
\end{equation}
It agrees exactly with \(\Phi\) under
\[
 (t,x,y_{\Phi},z_{\Phi})=(h,x,-z,y).
\]
The transformation \((y,z)\mapsto(-z,y)\) has determinant \(+1\), so the
two conventions give the same matrix, Lorentz form, and orientation.

Nothing in the \(\mathrm{SU}(3)\) orbit calculation selects the coefficient
\(1/2\).  More generally,
\begin{equation}
 Q_{4,\lambda}(r)=\lambda\left(\sum_ir_i\right)^2-\sum_ir_i^2
 =(4\lambda-1)h^2-x^2-y^2-z^2.
 \label{eq:hrgc-q4-lambda-family}
\end{equation}
Every \(\lambda>1/4\) therefore has Lorentz signature \((1,3)\), after a
rescaling of \(h\).  The value \(\lambda=1/2\) is the unique member for which
the fixed orthogonal coordinates in \eqref{eq:hrgc-hadamard-transform} have
equal unit coefficients.  This is a normalization choice in coefficient
space, not a theorem forced by the Gram quotient.

Nor does a nonnegative coefficient tuple automatically determine a point of
the positive Gram cone.  The tuple \((1,0,0,0)\) has
\(Q_4=-1/2\).  A coefficient tuple enters the Gram quotient through
\eqref{eq:hrgc-pro-hermitian-convention} only when that Hermitian matrix is
positive semidefinite, equivalently when \(h\geq0\) and \(Q_4\geq0\).
The hypersurface \(z=0\) is a real-symmetric coefficient slice, but that
fact alone does not supply a coordinatewise nonnegative embedding of a
three-component mass domain.

\subsection{The compact representation quotient}

\begin{theorem}[The \(S^{12}\) representation quotient]
\label{thm:hrgc-s12-ball-quotient}
Let
\[
 S^{12}_{\mathrm{col}}=\mathscr V^+=\mathscr V\cup\{\infty\}
\]
be the one-point compactification, and extend the \(\mathrm{SU}(3)\) action
by fixing \(\infty\).  Then
\begin{equation}
 S^{12}_{\mathrm{col}}/\mathrm{SU}(3)
 \cong
 \bigl(\operatorname{Herm}_2(\mathbb C)_{\geq0}\bigr)^+
 \cong B^4.
 \label{eq:hrgc-s12-ball}
\end{equation}
The orbit-type strata are
\begin{equation}
 \begin{array}{c|c|c}
 \text{source stratum}&\text{quotient stratum}&\text{stabilizer}\\ \hline
 \operatorname{rank}A=2&\operatorname{int}B^4&\{1\}\\
 \operatorname{rank}A=1&\partial B^4\setminus\{p_0,p_\infty\}&
   \mathrm{SU}(2)\\
 A=0&\{p_0\}&\mathrm{SU}(3)\\
 A=\infty&\{p_\infty\}&\mathrm{SU}(3).
 \end{array}
 \label{eq:hrgc-compact-strata}
\end{equation}
The ball identification is topological on the full stratified quotient.  On
its smooth interior, it may be chosen to carry the quotient orientation to
the orientation fixed in \eqref{eq:hrgc-pauli-coordinates}.
\end{theorem}

\begin{proof}
As a real vector space, \(\mathscr V\cong\mathbb R^{12}\), so its one-point
compactification is \(S^{12}\).  On
\(\mathscr V\cong\mathbb C^3\oplus\mathbb C^3\), the action of
\(U\in\mathrm{SU}(3)\) has complex determinant \((\det U)^2=1\), hence
real determinant one; it preserves the standard orientation before and after
one-point compactification.

The properness proved from \eqref{eq:hrgc-properness-identity} implies that
\(\Gamma\) extends continuously to
\[
 \Gamma^+:\mathscr V^+\longrightarrow
 \bigl(\operatorname{Herm}_2(\mathbb C)_{\geq0}\bigr)^+,
 \qquad \Gamma^+(\infty)=\infty.
\]
It separates the finite orbits by \cref{thm:hrgc-gram-quotient}, and the
added point is a distinct fixed orbit.  Hence it induces a continuous
bijection
\[
 \mathscr V^+/\mathrm{SU}(3)
 \longrightarrow
 \bigl(\operatorname{Herm}_2(\mathbb C)_{\geq0}\bigr)^+.
\]
The source is compact Hausdorff because a compact group acts on the compact
Hausdorff space \(\mathscr V^+\).  The target is Hausdorff because the
positive-semidefinite cone is locally compact Hausdorff.  The continuous
bijection is therefore a homeomorphism.

By \eqref{eq:hrgc-positive-cone-inequalities}, the map
\[
 (t,\mathbf x)\longmapsto(t-\lVert\mathbf x\rVert,\mathbf x)
\]
is a homeomorphism from the closed future cone to the closed half-space
\(\mathbb R_{\geq0}\times\mathbb R^3\).  For \(u=(s,\mathbf x)\) in that half-space,
inverse stereographic projection is
\begin{equation}
 u\longmapsto
 \frac{1}{1+\lVert u\rVert^2}
 \bigl(2s,2\mathbf x,\lVert u\rVert^2-1\bigr)
 \in S^4\subset\mathbb R^5.
 \label{eq:hrgc-halfspace-stereographic}
\end{equation}
Its image is the closed hemisphere with first coordinate nonnegative, minus
the north pole.  The point at infinity supplies that pole.  A closed
hemisphere is a closed four-ball, proving \eqref{eq:hrgc-s12-ball}.

Positive-definite matrices correspond to the open half-space and hence the
ball interior.  Nonzero rank-one matrices correspond to
\(\{0\}\times(\mathbb R^3\setminus\{0\})\).  The zero matrix and the point at
infinity give two distinct endpoints in the compactification of that
three-dimensional boundary.  Combining this with
\cref{prop:hrgc-stabilizers} proves \eqref{eq:hrgc-compact-strata}.  The
orientation statement on the connected interior is the calculation following
\eqref{eq:hrgc-orientation-determinant}.
\end{proof}

The notation \(S^{12}_{\mathrm{col}}\) records the construction.  There is a
canonical pointed homeomorphism of representation spheres
\[
 (\mathbb C^3)^+\wedge(\mathbb C^3)^+
 \cong(\mathbb C^3\oplus\mathbb C^3)^+
 =S^{12}_{\mathrm{col}}.
\]
If a different based six-sphere \(X\) is used, a based homeomorphism
\(X\cong(\mathbb C^3)^+\) transports the action to \(X\wedge X\).  The
transported action depends on that choice.  Without an independently proved
equivariance statement, it is not induced by a torus fibration, cusp
monodromy, period matrix, or fibrewise flux on \(X\).

\subsection{Exact gauge-theoretic boundary}

\begin{theorem}[No native nonzero \(S^4\) sector from the Gram quotient]
\label{thm:hrgc-no-native-s4-sector}
The construction above supplies no native principal \(\mathrm{SU}(3)\)
bundle of nonzero second Chern number on \(S^4\).  More precisely:
\begin{enumerate}
\item the orbit projection over the rank-two locus is the product bundle in
\eqref{eq:hrgc-free-product}, so its second Chern class is zero;
\item on the rank-one boundary the orbit projection is not a principal
\(\mathrm{SU}(3)\)-bundle, because its stabilizer is \(\mathrm{SU}(2)\);
\item the full compact quotient is \(B^4\), not \(S^4\); and
\item replacing \(B^4\) by the one-point compactification of its interior is
a different construction whose clutching class is additional data.
\end{enumerate}
\end{theorem}

\begin{proof}
The first two assertions are
\cref{thm:hrgc-free-trivialization,prop:hrgc-stabilizers}; the third is
\cref{thm:hrgc-s12-ball-quotient}.  For the fourth, the matrix logarithm is a
diffeomorphism
\[
 \log:\operatorname{Herm}_2(\mathbb C)_{>0}
 \xrightarrow{\ \cong\ }
 \operatorname{Herm}_2(\mathbb C)\cong\mathbb R^4.
\]
Its one-point compactification is therefore \(S^4\), but this operation
collapses the entire finite rank-one boundary of the full Gram quotient into
the new point at infinity.  Principal \(\mathrm{SU}(3)\)-bundles on this
new \(S^4\) are obtained by clutching two trivial bundles over hemispheres
with a map
\[
 \kappa:S^3\longrightarrow\mathrm{SU}(3).
\]
The fibration
\(\mathrm{SU}(2)\to\mathrm{SU}(3)\to S^5\), together with
\(\pi_4(S^5)=\pi_3(S^5)=0\), shows that
\[
 \pi_3(\mathrm{SU}(3))\cong\pi_3(\mathrm{SU}(2))
 =\pi_3(S^3)\cong\mathbb Z.
\]
Choose the generator so that its associated bundle has second Chern number
\(+1\).  A clutching class \(k\in\mathbb Z\) then has second Chern number
\(k\).  Extending the explicit product framing in
\cref{thm:hrgc-free-trivialization} by the same framing gives the constant
clutching map and \(k=0\).  Neither \(A^\dagger A\) nor
\(\lVert a\wedge b\rVert^2\) selects a nonconstant \(\kappa\).  Thus
\(k\ne0\) is valid additional gauge data, but it is not produced by the Gram
construction.
\end{proof}

There are two further cohomological checks.  If a candidate total space is
diffeomorphic to \(S^6\), then
\[
 H^2(S^6;\mathbb Z)=H^4(S^6;\mathbb Z)=0.
\]
Consequently no nonzero fibrewise degree-two class can be the restriction of
a degree-two class on that total space, and no \(c_2\) on the total space can
restrict to a nonzero degree-four class on a fibre.  This does not forbid
independently chosen higher characteristic data; it proves exactly that the
degree-two and degree-four classes in question do not extend.

Likewise, for every continuous map \(f:S^{12}\to S^4\) and generator
\(u_4\in H^4(S^4;\mathbb Z)\),
\[
 f^*u_4=0
\]
because \(H^4(S^{12};\mathbb Z)=0\).  Therefore the pullback along any such
map of an \(\mathrm{SU}(3)\) bundle with nonzero \(c_2=ku_4\) has
\(c_2=0\).  This statement does not assert that every intrinsic
\(\mathrm{SU}(3)\) bundle on \(S^{12}\) is trivial; it says that an
\(S^4\) second-Chern charge is not inherited by pullback to the
representation sphere.

\subsection{The log-Euclidean metric and Yang--Mills theory are extra choices}

The logarithm identifies the open positive cone with a real vector space.
One may therefore declare
\begin{equation}
 X=\log H,
 \qquad
 g_{\log}=\frac12\operatorname{tr}(dX\,dX),
 \label{eq:hrgc-log-metric}
\end{equation}
which is a complete flat Euclidean metric in the \(X\)-coordinates.  This
metric is not induced by the Frobenius metric on \(\mathscr V\).  To see the
difference without discarding any scale factor, set
\[
 H(s)=sI_2,
 \qquad
 A(s)=\sqrt{s}\begin{pmatrix}1&0\\0&1\\0&0\end{pmatrix},
 \qquad s>0.
\]
For \(\xi\in\mathfrak{su}(3)\),
\[
 \operatorname{Re}\operatorname{tr}
 \bigl(\dot A^\dagger\xi A\bigr)
 =\frac{\dot s}{2}
 \operatorname{Re}(\xi_{11}+\xi_{22})=0,
\]
so \(\dot A\) is horizontal for the Frobenius quotient.  Hence
\begin{equation}
 \lVert\dot H\rVert_{\mathrm{quot}}^2
 =\lVert\dot A\rVert_F^2=\frac{\dot s^2}{2s}.
 \label{eq:hrgc-frobenius-ray}
\end{equation}
By contrast, \(X(s)=(\log s)I_2\), and
\begin{equation}
 \lVert\dot H\rVert_{\log}^2
 =\frac12\operatorname{tr}(\dot X^2)
 =\frac{\dot s^2}{s^2}.
 \label{eq:hrgc-log-ray}
\end{equation}
The ratio of \eqref{eq:hrgc-frobenius-ray} and
\eqref{eq:hrgc-log-ray} depends on \(s\), so no global constant rescaling
identifies the metrics.  The Lorentz form
\eqref{eq:hrgc-lorentz-determinant} is a third object and is indefinite.

After choosing \eqref{eq:hrgc-log-metric}, one may define classical
Yang--Mills theory on the trivial product bundle supplied by
\cref{thm:hrgc-free-trivialization}.  The additional choice of the zero
connection has \(F_0=0\), so it solves the nonlinear classical equation
\(d_0^*F_0=0\); its Coulomb Hessian is the transverse free Laplacian.  These
are the standard zero-charge vacuum statements on a chosen Euclidean
\(\mathbb R^4\).  The Gram morphism itself supplies no nonzero clutching
class, no nonflat connection, no interacting Hamiltonian, no quantum
Hilbert-space identification, and no mass-gap estimate.  Thus the exact
logical implication is
\[
 \boxed{
 \begin{array}{c}
 \text{four-dimensional Gram/Lorentz quotient: proved;}\\
 S^{12}_{\mathrm{col}}/\mathrm{SU}(3)\cong B^4: \text{ proved;}\\
 \text{nonzero native }\mathrm{SU}(3)\text{ sector on }S^4:
 \text{ not produced;}\\
 \text{interacting or quantum Yang--Mills inheritance: not produced.}
 \end{array}}
\]
